\documentclass[11pt]{article}
\usepackage[margin=1in]{geometry}
\usepackage{amsmath,amssymb,amsthm,hyperref,enumitem}
\hypersetup{colorlinks=true,linkcolor=blue,urlcolor=blue}
\newtheorem{theorem}{Theorem}
\newtheorem{corollary}{Corollary}
\newcommand{\R}{\mathbb R}
\title{Two multi-affine equations with exponentially many\\real connected components}
\author{Principia Math In Progress}
\date{22 September 2026}
\begin{document}
\maketitle
\noindent\textbf{MathDB 405042. Negative answer to the degree-only bound.}

\section*{Task 1: the original question and its answer}
The question in Basu's contribution to Oberwolfach Report 9/2025
[1, pp.~445--446] asks whether two real multi-affine equations of
bounded total degree have a number of real connected components
bounded independently of the number of variables.
Here multi-affine means degree at most one in each individual variable.
The answer is no, already for total degrees two and four.

\begin{theorem}
For every integer \(m\geq2\), there are two multi-affine polynomials
of total degrees two and four in \(2m\) real variables whose common
real zero set has exactly \(2^m\) connected components.
Every component is homeomorphic to \(\R^m\).
\end{theorem}

\begin{proof}
Use independent variables \(x_1,y_1,\ldots,x_m,y_m\), and put
\[
 z_i=x_i y_i,\qquad
 F_m=\sum_{i=1}^m z_i-m,\qquad
 G_m=\sum_{1\leq i<j\leq m}z_i z_j-\binom m2.                 \tag{1}
\]
Every nonconstant monomial of \(F_m\) is \(x_i y_i\).
Every nonconstant monomial of \(G_m\) is
\(x_i y_i x_j y_j\), where \(i\ne j\).
Thus no variable occurs with exponent greater than one in either
polynomial. Their total degrees are exactly two and four.

The polynomial identity
\[
 \sum_{i=1}^m(z_i-1)^2
       =F_m^2+2(m-1)F_m-2G_m                              \tag{2}
\]
follows by expanding the square of \(\sum_i z_i\).
At a common real zero of \(F_m,G_m\), its left side vanishes,
so \(z_i=1\) for every \(i\). Conversely, these equalities imply
\(F_m=G_m=0\). Consequently the common real zero set is exactly
\[
 V_m=\{(x,y):x_i y_i=1\text{ for every }i\}
       \cong(\R\setminus\{0\})^m.                         \tag{3}
\]
The homeomorphism sends \((x_1,\ldots,x_m)\) to
\((x_1,x_1^{-1},\ldots,x_m,x_m^{-1})\); its inverse is the
projection to the \(x\)-coordinates.

Each \(x_i\) is nonzero on \(V_m\). The map taking a point to the
sign vector of its \(x\)-coordinates is continuous into the discrete
set \(\{-1,1\}^m\). Hence a connected subset has a constant sign vector.
For each \(\epsilon\in\{-1,1\}^m\), the corresponding fiber is
homeomorphic to \(\R^m\) by
\[
 t\longmapsto
  (\epsilon_1e^{t_1},\epsilon_1e^{-t_1},
        \ldots,\epsilon_me^{t_m},\epsilon_me^{-t_m}).       \tag{4}
\]
Its inverse is \(t_i=\log|x_i|\). Thus all \(2^m\) fibers are nonempty
and connected, and they are precisely the connected components.
\end{proof}

The degrees stay bounded by four while \(m\) is arbitrary, so (3)
rules out any finite bound depending only on those degrees.
Only the subsequence of even ambient dimensions is needed.
An unused free variable gives the same component count in dimension
\(2m+1\). The coefficients are allowed to depend on \(m\), as in the
source question. Identity (2) is a proof device: its square is not
being introduced as one of the two defining equations.

\section*{A smooth strengthening}
The conclusion does not depend on the singularity of the pair (1).
Fix \(0<\rho<1\) and replace its second equation by
\[
 G_{m,\rho}=G_m+\rho/2=0.                                 \tag{5}
\]
\begin{corollary}
For every \(m\geq3\), the two equations \(F_m=G_{m,\rho}=0\)
define a smooth real submanifold of dimension \(2m-2\), with
exactly \(2^m\) connected components. Each component is
diffeomorphic to \(\R^m\times S^{m-2}\).
\end{corollary}
\begin{proof}
Identity (2) gives an exact description in the \(z\)-variables:
\[
 Z_{m,\rho}
 =\left\{z:\sum_i z_i=m,\ \sum_i(z_i-1)^2=\rho\right\}.     \tag{6}
\]
This is a sphere of radius \(\sqrt\rho\), centered at
\((1,\ldots,1)\), in the \((m-1)\)-dimensional affine hyperplane
\(\sum_i z_i=m\). In particular it is diffeomorphic to \(S^{m-2}\).
It is nonempty, and \(z_i\geq1-\sqrt\rho>0\) throughout.

Fix a sign vector \(\epsilon\). The map from
\(\R^m\times Z_{m,\rho}\) given by
\[
 x_i=\epsilon_i e^{t_i},\qquad
 y_i=\epsilon_i z_i e^{-t_i}                               \tag{7}
\]
is a diffeomorphism onto the part with these signs.
Its inverse is \(t_i=\log|x_i|\), \(z_i=x_i y_i\).
Since \(m\geq3\), the sphere in (6) is connected. Each sign fiber
is therefore connected; the continuous sign map again separates
the \(2^m\) fibers.

For completeness, the two defining equations are transverse at
every real solution. The map \((x,y)\mapsto z\) has surjective
differential because \(x_i\ne0\), and its derivatives with respect
to the \(y_i\) have nonzero diagonal entries \(x_i\).
In \(z\)-coordinates the two gradients are
\[
 (1,\ldots,1),\qquad
 \left(\sum_jz_j-z_1,\ldots,\sum_jz_j-z_m\right).
\]
They can be dependent only if all \(z_i\) coincide. On the
hyperplane in (6), that would give \(z_i=1\), contradicting
\(\rho>0\). Their pullbacks are independent, proving the claimed
smoothness and dimension.
\end{proof}

\section*{Tasks 2--5: four proposed mathematical follow-ups}
\begin{enumerate}[label=\textbf{Task \arabic*:},start=2,leftmargin=*]
\item Determine the least total-degree bound \(d\) for which two
multi-affine equations can have unboundedly many connected
components as the ambient dimension grows; in particular resolve
the cases \(d=2\) and \(d=3\).
\item Determine the largest possible number of real connected
components of two multi-affine equations of degrees at most four
in \(n\) variables, or determine its exponential growth rate.
\item Decide whether a degree-only component bound holds when
each of the two defining multi-affine polynomials is invariant
under every permutation of the \(n\) variables.
\item Determine whether two multi-affine equations of uniformly
bounded total degree can have a compact smooth common real zero
set with an unbounded number of connected components.
\end{enumerate}
These are proposed directions, not additional claimed solutions.

\section*{Verification, sources, and scope of review}
The accompanying exact verifier checks the identity, monomial
degrees, representatives of every sign component in finite
dimensions, and rational transverse examples of (5). These checks
supplement the proofs for arbitrary \(m\); they do not prove the
infinite family. Review is by the authoring agent only, without
independent peer review or a proof assistant.

\medskip\noindent
[1] \emph{MFO--RIMS Tandem Workshop: Optimization, Theoretical
Computer Science and Algebraic Geometry: Convexity and Beyond},
Oberwolfach Report 9/2025,
\href{https://doi.org/10.4171/OWR/2025/9}{doi:10.4171/OWR/2025/9}.
The complete 46-page published report was read. The contribution
on printed pp.~445--446 states the unrestricted two-equation
question and cites the earlier three-equation construction.

\medskip\noindent
[2] S. Basu and D. Perrucci,
\emph{Topology of real multi-affine hypersurfaces and a
homological stability property},
\href{https://arxiv.org/abs/2204.01595}{arXiv:2204.01595v1}
(4 April 2022), all 26 pages read.
Its Example 2.2 uses three equations; its two-equation question
is open on pp.~6 and 24.
The related 33-page journal version, \emph{Adv. Math.} 420 (2023),
108982, was identified but not obtained.
The later published source [1] explicitly cites that journal
paper and still poses the present question.

\medskip\noindent
The dated audit records the live MathDB checks and bounded
public searches. No earlier full solution was found in the
checked material; the audit is not a guarantee of exhaustive
priority coverage. The smooth strengthening receives no separate
campaign credit.
\end{document}
